\section{Conclusion}

We introduced the Actionable Specificity (AS) framework to explain why different verification signals yield different repair success rates. Our existence test revealed that regex-based AS extraction achieves only 49.5\% coverage, with near-zero state extraction (2.7\%). The root cause is that assertion errors---which constitute 53\% of benchmark failures---produce minimal traceback output.

\textbf{We acknowledge this as a limitation of our operationalization}, not evidence that the AS concept is invalid. The ordering preservation (C1$\geq$C2$\geq$C3$\geq$C4) on extractable signals suggests the decomposition may remain useful under alternative operationalizations.

\textbf{Contributions:}
\begin{enumerate}
    \item The AS framework provides a principled decomposition for analyzing feedback informativeness.
    \item We identified a necessary precondition: error type stratification before extraction.
    \item The framework provides actionable guidance: it identifies which error types block extraction and why.
\end{enumerate}

\textbf{Future Work:} Testing AS on traceback-rich error subsets, LLM-based extraction, and \texttt{pytest --tb=long} for richer diagnostics.
